% assets/w07-pc-fan.svg — p_c 를 바꾸면 그 자리의 화살표가 어떻게 돌아가는가
%
%   bash _tools/tikz2svg.sh assets/src/w07-pc-fan.tex assets/w07-pc-fan.svg
%
% 요점: 같은 자리(r = 32 m, r_d = 25 m)에서 p_c 만 바꿔 본다.
%       부호는 화살표를 e_theta 의 어느 쪽에 놓을지를, 크기는 e_r 반대쪽에
%       얼마나 가깝게 놓을지를 정한다.
%
%   alpha = atan2(p_c * r, -(r - r_d)) = atan2(p_c*32, -7)
%     p_c =  0.150 -> atan2( 4.80, -7) =  145.6 deg
%     p_c =  0.313 -> atan2(10.02, -7) =  124.9 deg
%     p_c =  1.000 -> atan2(32.00, -7) =  102.3 deg
%     p_c = -0.313 -> atan2(-10.02,-7) = -124.9 deg
%     p_c = -1.000 -> atan2(-32.00,-7) = -102.3 deg
%   TikZ 각 = (90 - theta) - alpha = 40 - alpha

\documentclass[border=6pt]{standalone}
% capstone-style.tex — 이 볼트의 TikZ 그림이 공유하는 스타일
%
% 쓰는 법:  % capstone-style.tex — 이 볼트의 TikZ 그림이 공유하는 스타일
%
% 쓰는 법:  % capstone-style.tex — 이 볼트의 TikZ 그림이 공유하는 스타일
%
% 쓰는 법:  \input{capstone-style}   (assets/src/ 안에서)
% 빌드:     bash _tools/tikz2svg.sh assets/src/w02-node-topic.tex assets/w02-node-topic.svg
%
% 왜 TikZ 인가
%   손으로 SVG 좌표를 쓰면 화살촉·선 굵기·글자 크기가 그림마다 미묘하게 달라진다.
%   그것이 "자료가 어설퍼 보인다"의 정체다. 스타일을 한 번만 정하고 상대좌표로
%   배치하면 좌표를 고쳐도 그림이 무너지지 않는다.
%   수식은 본문과 같은 TeX 글꼴로 나온다.
%
% 한글
%   ko.TeX(DVI 경로)를 쓴다. XeLaTeX 은 TikZ 그래픽을 PDF special 로 내보내는데
%   dvisvgm 이 그것을 받으려면 Ghostscript < 10.01 이 필요하다. 이 PC 에는
%   10.07 이 깔려 있어 그 경로가 막힌다. latex -> DVI -> dvisvgm 이 유일하게
%   동작하는 조합이다.

\usepackage[hangul]{kotex}
\usepackage{tikz}
\usepackage{amsmath}
\usepackage{xcolor}
\usetikzlibrary{arrows.meta, positioning, calc, fit, backgrounds, decorations.pathmorphing}

% ---- 색 --------------------------------------------------------------------
% 강조는 하나만. 나머지는 검정. 흑백 인쇄에서도 읽혀야 한다.
\definecolor{ink}{HTML}{1A1A1A}
\definecolor{soft}{HTML}{666666}
\definecolor{accent}{HTML}{7C3AED}
\definecolor{warn}{HTML}{B91C1C}

% 계층 색 — Simulink 모델의 블록 색과 같은 규칙을 쓴다
\definecolor{guid}{HTML}{1D4ED8}    % 유도
\definecolor{ctrl}{HTML}{EA580C}    % 제어
\definecolor{thr}{HTML}{D97706}     % 추진기
\definecolor{plant}{HTML}{16A34A}   % 운동모델
\definecolor{nav}{HTML}{7C3AED}     % 항법 · 신호처리
\definecolor{pathc}{HTML}{0D9488}   % 계획 경로 · 항적
\definecolor{errc}{HTML}{C2410C}    % 오차
\definecolor{flow}{HTML}{0369A1}    % 조류 · 바람

% ---- 선과 화살촉 -----------------------------------------------------------
\tikzset{
  sig/.style   = {ink, line width=0.5pt, -{Stealth[length=4pt, width=3pt]}},
  wire/.style  = {ink, line width=0.5pt},
  fb/.style    = {ink, line width=0.5pt, densely dashed,
                  -{Stealth[length=4pt, width=3pt]}},
  asig/.style  = {accent, line width=0.5pt, densely dashed,
                  -{Stealth[length=4pt, width=3pt]}},
  % 블록 — 흰 바탕, 직각, 검은 테두리. 색은 테두리에만 준다
  blk/.style   = {draw=ink, line width=0.55pt, fill=white,
                  minimum width=18mm, minimum height=8.5mm, inner sep=2pt, font=\small},
  gblk/.style  = {blk, draw=guid},
  cblk/.style  = {blk, draw=ctrl},
  tblk/.style  = {blk, draw=thr},
  pblk/.style  = {blk, draw=plant},
  nblk/.style  = {blk, draw=nav},
  % 합산점 · 분기점
  sum/.style   = {draw=ink, line width=0.55pt, fill=white, circle,
                  inner sep=0pt, minimum size=4.6mm, font=\scriptsize},
  dot/.style   = {ink, fill=ink, circle, inner sep=0pt, minimum size=1.5mm},
  % 글자
  sname/.style = {font=\small, inner sep=1.5pt},
  note/.style  = {font=\scriptsize, soft, inner sep=1.5pt},
  wnote/.style = {font=\scriptsize, warn, inner sep=1.5pt},
  anote/.style = {font=\scriptsize, accent, inner sep=1.5pt},
  axis/.style  = {ink, line width=0.5pt},
  tick/.style  = {font=\scriptsize, soft},
  % 묶음 상자 — 서브시스템 표시
  grp/.style   = {draw=soft, line width=0.4pt, densely dotted, rounded corners=1mm,
                  inner sep=3mm},
}
   (assets/src/ 안에서)
% 빌드:     bash _tools/tikz2svg.sh assets/src/w02-node-topic.tex assets/w02-node-topic.svg
%
% 왜 TikZ 인가
%   손으로 SVG 좌표를 쓰면 화살촉·선 굵기·글자 크기가 그림마다 미묘하게 달라진다.
%   그것이 "자료가 어설퍼 보인다"의 정체다. 스타일을 한 번만 정하고 상대좌표로
%   배치하면 좌표를 고쳐도 그림이 무너지지 않는다.
%   수식은 본문과 같은 TeX 글꼴로 나온다.
%
% 한글
%   ko.TeX(DVI 경로)를 쓴다. XeLaTeX 은 TikZ 그래픽을 PDF special 로 내보내는데
%   dvisvgm 이 그것을 받으려면 Ghostscript < 10.01 이 필요하다. 이 PC 에는
%   10.07 이 깔려 있어 그 경로가 막힌다. latex -> DVI -> dvisvgm 이 유일하게
%   동작하는 조합이다.

\usepackage[hangul]{kotex}
\usepackage{tikz}
\usepackage{amsmath}
\usepackage{xcolor}
\usetikzlibrary{arrows.meta, positioning, calc, fit, backgrounds, decorations.pathmorphing}

% ---- 색 --------------------------------------------------------------------
% 강조는 하나만. 나머지는 검정. 흑백 인쇄에서도 읽혀야 한다.
\definecolor{ink}{HTML}{1A1A1A}
\definecolor{soft}{HTML}{666666}
\definecolor{accent}{HTML}{7C3AED}
\definecolor{warn}{HTML}{B91C1C}

% 계층 색 — Simulink 모델의 블록 색과 같은 규칙을 쓴다
\definecolor{guid}{HTML}{1D4ED8}    % 유도
\definecolor{ctrl}{HTML}{EA580C}    % 제어
\definecolor{thr}{HTML}{D97706}     % 추진기
\definecolor{plant}{HTML}{16A34A}   % 운동모델
\definecolor{nav}{HTML}{7C3AED}     % 항법 · 신호처리
\definecolor{pathc}{HTML}{0D9488}   % 계획 경로 · 항적
\definecolor{errc}{HTML}{C2410C}    % 오차
\definecolor{flow}{HTML}{0369A1}    % 조류 · 바람

% ---- 선과 화살촉 -----------------------------------------------------------
\tikzset{
  sig/.style   = {ink, line width=0.5pt, -{Stealth[length=4pt, width=3pt]}},
  wire/.style  = {ink, line width=0.5pt},
  fb/.style    = {ink, line width=0.5pt, densely dashed,
                  -{Stealth[length=4pt, width=3pt]}},
  asig/.style  = {accent, line width=0.5pt, densely dashed,
                  -{Stealth[length=4pt, width=3pt]}},
  % 블록 — 흰 바탕, 직각, 검은 테두리. 색은 테두리에만 준다
  blk/.style   = {draw=ink, line width=0.55pt, fill=white,
                  minimum width=18mm, minimum height=8.5mm, inner sep=2pt, font=\small},
  gblk/.style  = {blk, draw=guid},
  cblk/.style  = {blk, draw=ctrl},
  tblk/.style  = {blk, draw=thr},
  pblk/.style  = {blk, draw=plant},
  nblk/.style  = {blk, draw=nav},
  % 합산점 · 분기점
  sum/.style   = {draw=ink, line width=0.55pt, fill=white, circle,
                  inner sep=0pt, minimum size=4.6mm, font=\scriptsize},
  dot/.style   = {ink, fill=ink, circle, inner sep=0pt, minimum size=1.5mm},
  % 글자
  sname/.style = {font=\small, inner sep=1.5pt},
  note/.style  = {font=\scriptsize, soft, inner sep=1.5pt},
  wnote/.style = {font=\scriptsize, warn, inner sep=1.5pt},
  anote/.style = {font=\scriptsize, accent, inner sep=1.5pt},
  axis/.style  = {ink, line width=0.5pt},
  tick/.style  = {font=\scriptsize, soft},
  % 묶음 상자 — 서브시스템 표시
  grp/.style   = {draw=soft, line width=0.4pt, densely dotted, rounded corners=1mm,
                  inner sep=3mm},
}
   (assets/src/ 안에서)
% 빌드:     bash _tools/tikz2svg.sh assets/src/w02-node-topic.tex assets/w02-node-topic.svg
%
% 왜 TikZ 인가
%   손으로 SVG 좌표를 쓰면 화살촉·선 굵기·글자 크기가 그림마다 미묘하게 달라진다.
%   그것이 "자료가 어설퍼 보인다"의 정체다. 스타일을 한 번만 정하고 상대좌표로
%   배치하면 좌표를 고쳐도 그림이 무너지지 않는다.
%   수식은 본문과 같은 TeX 글꼴로 나온다.
%
% 한글
%   ko.TeX(DVI 경로)를 쓴다. XeLaTeX 은 TikZ 그래픽을 PDF special 로 내보내는데
%   dvisvgm 이 그것을 받으려면 Ghostscript < 10.01 이 필요하다. 이 PC 에는
%   10.07 이 깔려 있어 그 경로가 막힌다. latex -> DVI -> dvisvgm 이 유일하게
%   동작하는 조합이다.

\usepackage[hangul]{kotex}
\usepackage{tikz}
\usepackage{amsmath}
\usepackage{xcolor}
\usetikzlibrary{arrows.meta, positioning, calc, fit, backgrounds, decorations.pathmorphing}

% ---- 색 --------------------------------------------------------------------
% 강조는 하나만. 나머지는 검정. 흑백 인쇄에서도 읽혀야 한다.
\definecolor{ink}{HTML}{1A1A1A}
\definecolor{soft}{HTML}{666666}
\definecolor{accent}{HTML}{7C3AED}
\definecolor{warn}{HTML}{B91C1C}

% 계층 색 — Simulink 모델의 블록 색과 같은 규칙을 쓴다
\definecolor{guid}{HTML}{1D4ED8}    % 유도
\definecolor{ctrl}{HTML}{EA580C}    % 제어
\definecolor{thr}{HTML}{D97706}     % 추진기
\definecolor{plant}{HTML}{16A34A}   % 운동모델
\definecolor{nav}{HTML}{7C3AED}     % 항법 · 신호처리
\definecolor{pathc}{HTML}{0D9488}   % 계획 경로 · 항적
\definecolor{errc}{HTML}{C2410C}    % 오차
\definecolor{flow}{HTML}{0369A1}    % 조류 · 바람

% ---- 선과 화살촉 -----------------------------------------------------------
\tikzset{
  sig/.style   = {ink, line width=0.5pt, -{Stealth[length=4pt, width=3pt]}},
  wire/.style  = {ink, line width=0.5pt},
  fb/.style    = {ink, line width=0.5pt, densely dashed,
                  -{Stealth[length=4pt, width=3pt]}},
  asig/.style  = {accent, line width=0.5pt, densely dashed,
                  -{Stealth[length=4pt, width=3pt]}},
  % 블록 — 흰 바탕, 직각, 검은 테두리. 색은 테두리에만 준다
  blk/.style   = {draw=ink, line width=0.55pt, fill=white,
                  minimum width=18mm, minimum height=8.5mm, inner sep=2pt, font=\small},
  gblk/.style  = {blk, draw=guid},
  cblk/.style  = {blk, draw=ctrl},
  tblk/.style  = {blk, draw=thr},
  pblk/.style  = {blk, draw=plant},
  nblk/.style  = {blk, draw=nav},
  % 합산점 · 분기점
  sum/.style   = {draw=ink, line width=0.55pt, fill=white, circle,
                  inner sep=0pt, minimum size=4.6mm, font=\scriptsize},
  dot/.style   = {ink, fill=ink, circle, inner sep=0pt, minimum size=1.5mm},
  % 글자
  sname/.style = {font=\small, inner sep=1.5pt},
  note/.style  = {font=\scriptsize, soft, inner sep=1.5pt},
  wnote/.style = {font=\scriptsize, warn, inner sep=1.5pt},
  anote/.style = {font=\scriptsize, accent, inner sep=1.5pt},
  axis/.style  = {ink, line width=0.5pt},
  tick/.style  = {font=\scriptsize, soft},
  % 묶음 상자 — 서브시스템 표시
  grp/.style   = {draw=soft, line width=0.4pt, densely dotted, rounded corners=1mm,
                  inner sep=3mm},
}


\begin{document}
\begin{tikzpicture}[x=1mm, y=1mm]

\def\Rd{32.5}
\def\Rp{41.6}
\def\Aer{40}
\def\Aet{-50}

\coordinate (O) at (0,0);
\coordinate (P) at (40:\Rp mm);

% ---- 목표 원 · 중심 -------------------------------------------------------
\draw[pathc, line width=2.0pt, opacity=0.16] (O) circle (\Rd mm);
\fill[ink] (O) circle (1.2);
\node[note, fill=white, inner sep=0.6pt, anchor=north] at (2.0,-2.0) {중심};
\draw[pathc, line width=0.7pt] (O) -- (215:\Rd mm);
\node[note, pathc, fill=white, inner sep=0.8pt] at (215:19mm) {$r_d = 25$\,m};
\draw[ink, line width=0.6pt, densely dotted] (O) -- (P);
\node[note, fill=white, inner sep=0.8pt] at (40:27mm) {$r = 32$\,m};

% ---- 배 자리의 두 기준 방향 -----------------------------------------------
\draw[soft, line width=0.4pt, densely dashed, -{Stealth[length=3pt,width=2.4pt]}]
      (P) -- ++(\Aer:17mm);
\node[note, anchor=south west, inner sep=1pt] at ($(P)+(\Aer:17.5mm)$) {$\hat{e}_r$};
\draw[soft, line width=0.4pt, densely dashed, -{Stealth[length=3pt,width=2.4pt]}]
      (P) -- ++(\Aet:30mm);
\node[note, anchor=west, inner sep=1pt] at ($(P)+(\Aet:31mm)$) {$\hat{e}_\theta$};

% ---- p_c > 0 — 시계 방향 ---------------------------------------------------
\foreach \pc/\dir/\lab in {0.15/-105.6/{$p_c = 0.15$}, 0.313/-84.9/{$p_c = 0.313$}, 1.0/-62.3/{$p_c = 1.0$}}{
  \draw[guid, line width=1.0pt, -{Stealth[length=4.5pt,width=3.5pt]}]
       (P) -- ++(\dir:26mm);
}
\node[note, guid, anchor=north east, inner sep=1pt] at ($(P)+(-105.6:27mm)$) {$p_c = 0.15$};
\node[note, guid, anchor=north,      inner sep=1pt] at ($(P)+(-84.9:27.5mm)$) {$0.313$};
\node[note, guid, anchor=north west, inner sep=1pt] at ($(P)+(-62.3:27mm)$) {$1.0$};

% ---- p_c < 0 — 반시계 방향 -------------------------------------------------
\foreach \dir in {164.9, 142.3}{
  \draw[errc, line width=1.0pt, -{Stealth[length=4.5pt,width=3.5pt]}]
       (P) -- ++(\dir:26mm);
}
\node[note, errc, anchor=east,       inner sep=1pt] at ($(P)+(164.9:27mm)$) {$p_c = -0.313$};
\node[note, errc, anchor=south east, inner sep=1pt] at ($(P)+(142.3:27mm)$) {$-1.0$};

% ---- 배 ---------------------------------------------------------------------
\fill[ink] (P) circle (1.6);
\node[sname, anchor=south, fill=white, inner sep=1pt] at ($(P)+(0,2.2)$) {배};

% ---- alpha 호 두 개 (양·음 각각 하나만) ------------------------------------
\draw[guid, line width=0.5pt, densely dotted, -{Stealth[length=3pt,width=2.4pt]}]
      ($(P)+(\Aer:11mm)$) arc (\Aer:-84.9:11mm);
\node[note, guid, fill=white, inner sep=0.8pt] at ($(P)+(-22:13mm)$) {$\alpha>0$};
\draw[errc, line width=0.5pt, densely dotted, -{Stealth[length=3pt,width=2.4pt]}]
      ($(P)+(\Aer:8mm)$) arc (\Aer:164.9:8mm);
\node[note, errc, fill=white, inner sep=0.8pt] at ($(P)+(104:10mm)$) {$\alpha<0$};

% ---- 회전 방향 미니 그림 (왼쪽 아래) ---------------------------------------
\begin{scope}[shift={(-46,-46)}]
  \draw[guid, line width=1.2pt, -{Stealth[length=7pt,width=5pt]}]
        (120:11mm) arc (120:-170:11mm);
  \node[note, guid, anchor=north, text width=26mm, align=center] at (0,-13.5)
       {$p_c>0$\\ 시계 방향};
\end{scope}
\begin{scope}[shift={(-10,-46)}]
  \draw[errc, line width=1.2pt, -{Stealth[length=7pt,width=5pt]}]
        (60:11mm) arc (60:350:11mm);
  \node[note, errc, anchor=north, text width=26mm, align=center] at (0,-13.5)
       {$p_c<0$\\ 반시계 방향};
\end{scope}

% ---- 읽는 법 ---------------------------------------------------------------
\node[anchor=north west, text width=66mm] at (28,-34)
     {\small $\alpha = \operatorname{atan2}\!\big(p_c\,r,\; -(r-r_d)\big)$};

\node[note, anchor=north west, text width=68mm] at (28,-43)
     {\textbf{부호} — $\alpha$ 가 $\hat{e}_\theta$ 쪽(양)인지 반대쪽(음)인지만 바뀐다.
      거리를 좁히는 $\dot{r}$ 성분은 그대로다\\[2pt]
      \textbf{크기} — $|p_c|$ 가 작을수록 $\alpha$ 가 $180^\circ$ 에 가까워진다.
      $180^\circ$ 는 $\hat{e}_r$ 의 정반대 = 중심을 곧장 향함\\[2pt]
      $|p_c|$ 0.15 / 0.313 / 1.0 $\rightarrow$ $\alpha$ $145.6^\circ$ / $124.9^\circ$ / $102.3^\circ$\\[2pt]
      $|p_c|$ 가 커질수록 $\alpha \rightarrow 90^\circ$ — 거의 돌기만 하고 원으로 붙지 않는다};

\node[wnote, anchor=north west, text width=112mm] at (-50,-72)
     {$r$ 가 바뀌면 같은 $p_c$ 라도 $\alpha$ 가 바뀐다 — 멀리 있을수록 $-(r-r_d)$ 가 커져
      $\alpha$ 가 $180^\circ$ 쪽으로 간다. 그래서 먼 데서는 중심 쪽으로 곧장 오고,
      원에 가까워질수록 접선 쪽으로 눕는다. 이것이 한 식으로 접근과 선회를 같이 처리하는 방식이다};

\end{tikzpicture}
\end{document}
